\chapter{The Command Line}
\label{chap:command-line}

\chapterintro{The command line is one of the most important tools for engineers. It lets you navigate files, run programs, automate work, and control development environments.}

\learningobjectives{
\begin{itemize}
    \item Explain what a shell is.
    \item Use basic terminal commands.
    \item Understand standard input, standard output, and standard error.
    \item Run Python scripts from the command line.
    \item Connect command-line skills to Git, Databricks, and deployment workflows.
\end{itemize}}

\section{What Is the Command Line?}

The command line is a text-based interface for controlling a computer. Instead of clicking buttons, you type commands.

\begin{definition}
A shell is a program that reads commands from a user, asks the operating system to execute them, and displays the results.
\end{definition}

Common shells include Bash, Zsh, PowerShell, and Fish.

\section{Why Engineers Use the Command Line}

The command line is useful because it is:

\begin{itemize}
    \item fast,
    \item scriptable,
    \item precise,
    \item available on servers,
    \item essential for Git and deployment.
\end{itemize}

Many production systems do not have graphical interfaces. Engineers must be comfortable working with text commands.

\section{Basic Commands}

\begin{center}
\begin{tabular}{ll}
\toprule
Command & Meaning \\
\midrule
\texttt{pwd} & Print current directory \\
\texttt{ls} & List files \\
\texttt{cd} & Change directory \\
\texttt{mkdir} & Create directory \\
\texttt{touch} & Create empty file \\
\texttt{cat} & Print file contents \\
\texttt{cp} & Copy files \\
\texttt{mv} & Move or rename files \\
\texttt{rm} & Remove files \\
\bottomrule
\end{tabular}
\end{center}

\section{A First Terminal Session}

\begin{lstlisting}[language=bash]
pwd
mkdir command-line-lab
cd command-line-lab
touch notes.txt
echo "Hello from the terminal" > notes.txt
cat notes.txt
\end{lstlisting}

\section{Standard Streams}

Programs commonly use three streams.

\begin{itemize}
    \item Standard input: information entering the program.
    \item Standard output: normal program output.
    \item Standard error: error messages.
\end{itemize}

\section{Running Python from the Terminal}

Create a file named \texttt{main.py}.

\begin{lstlisting}
print("Hello from Python")
\end{lstlisting}

Run it:

\begin{lstlisting}[language=bash]
python main.py
\end{lstlisting}

\section{Why This Matters for LLM Engineering}

LLM engineers use command-line tools to:

\begin{itemize}
    \item install packages,
    \item run tests,
    \item start servers,
    \item manage Git repositories,
    \item launch training jobs,
    \item inspect logs,
    \item deploy applications.
\end{itemize}

\section{Git Checkpoint}

\begin{lstlisting}[language=bash]
git add volumes/volume01/chapters/chapter12.tex
git commit -m "feat(volume01): add command line chapter"
\end{lstlisting}

\section{Exercises}

\begin{exercise}
Use the command line to create a folder called \texttt{terminal-practice}, create a file inside it, write one line into the file, and print the file contents.
\end{exercise}

\begin{solution}
One possible solution is:
\begin{lstlisting}[language=bash]
mkdir terminal-practice
cd terminal-practice
echo "practice" > notes.txt
cat notes.txt
\end{lstlisting}
\end{solution}

\section{Portfolio Milestone}

Create a command-line lab repository containing a README, a Python script, and a transcript of at least ten terminal commands you successfully executed.

\section{Summary}

The command line is a text interface for controlling a computer. It is essential for professional engineering because it supports automation, remote work, Git, package management, testing, and deployment.
